\IfFileExists{stacks-project.cls}{%
\documentclass{stacks-project}
}{%
\documentclass{amsart}
}

% For dealing with references we use the comment environment
\usepackage{verbatim}
\newenvironment{reference}{\comment}{\endcomment}
%\newenvironment{reference}{}{}
\newenvironment{slogan}{\comment}{\endcomment}
\newenvironment{history}{\comment}{\endcomment}

% For commutative diagrams we use Xy-pic
\usepackage[all]{xy}

% We use 2cell for 2-commutative diagrams.
\xyoption{2cell}
\UseAllTwocells

% We use multicol for the list of chapters between chapters
\usepackage{multicol}

% This is generally recommended for better output
\usepackage{lmodern}
\usepackage[T1]{fontenc}

% For cross-file-references
\usepackage{xr-hyper}

% Package for hypertext links:
\usepackage{hyperref}

% For any local file, say "hello.tex" you want to link to please
% use \externaldocument[hello-]{hello}
\externaldocument[introduction-]{introduction}
\externaldocument[conventions-]{conventions}
\externaldocument[sets-]{sets}
\externaldocument[categories-]{categories}
\externaldocument[topology-]{topology}
\externaldocument[sheaves-]{sheaves}
\externaldocument[sites-]{sites}
\externaldocument[stacks-]{stacks}
\externaldocument[fields-]{fields}
\externaldocument[algebra-]{algebra}
\externaldocument[brauer-]{brauer}
\externaldocument[homology-]{homology}
\externaldocument[derived-]{derived}
\externaldocument[simplicial-]{simplicial}
\externaldocument[more-algebra-]{more-algebra}
\externaldocument[smoothing-]{smoothing}
\externaldocument[modules-]{modules}
\externaldocument[sites-modules-]{sites-modules}
\externaldocument[injectives-]{injectives}
\externaldocument[cohomology-]{cohomology}
\externaldocument[sites-cohomology-]{sites-cohomology}
\externaldocument[dga-]{dga}
\externaldocument[dpa-]{dpa}
\externaldocument[sdga-]{sdga}
\externaldocument[hypercovering-]{hypercovering}
\externaldocument[schemes-]{schemes}
\externaldocument[constructions-]{constructions}
\externaldocument[properties-]{properties}
\externaldocument[morphisms-]{morphisms}
\externaldocument[coherent-]{coherent}
\externaldocument[divisors-]{divisors}
\externaldocument[limits-]{limits}
\externaldocument[varieties-]{varieties}
\externaldocument[topologies-]{topologies}
\externaldocument[descent-]{descent}
\externaldocument[perfect-]{perfect}
\externaldocument[more-morphisms-]{more-morphisms}
\externaldocument[flat-]{flat}
\externaldocument[groupoids-]{groupoids}
\externaldocument[more-groupoids-]{more-groupoids}
\externaldocument[etale-]{etale}
\externaldocument[chow-]{chow}
\externaldocument[intersection-]{intersection}
\externaldocument[pic-]{pic}
\externaldocument[weil-]{weil}
\externaldocument[adequate-]{adequate}
\externaldocument[dualizing-]{dualizing}
\externaldocument[duality-]{duality}
\externaldocument[discriminant-]{discriminant}
\externaldocument[derham-]{derham}
\externaldocument[local-cohomology-]{local-cohomology}
\externaldocument[algebraization-]{algebraization}
\externaldocument[curves-]{curves}
\externaldocument[resolve-]{resolve}
\externaldocument[models-]{models}
\externaldocument[functors-]{functors}
\externaldocument[equiv-]{equiv}
\externaldocument[pione-]{pione}
\externaldocument[etale-cohomology-]{etale-cohomology}
\externaldocument[proetale-]{proetale}
\externaldocument[relative-cycles-]{relative-cycles}
\externaldocument[more-etale-]{more-etale}
\externaldocument[trace-]{trace}
\externaldocument[crystalline-]{crystalline}
\externaldocument[spaces-]{spaces}
\externaldocument[spaces-properties-]{spaces-properties}
\externaldocument[spaces-morphisms-]{spaces-morphisms}
\externaldocument[decent-spaces-]{decent-spaces}
\externaldocument[spaces-cohomology-]{spaces-cohomology}
\externaldocument[spaces-limits-]{spaces-limits}
\externaldocument[spaces-divisors-]{spaces-divisors}
\externaldocument[spaces-over-fields-]{spaces-over-fields}
\externaldocument[spaces-topologies-]{spaces-topologies}
\externaldocument[spaces-descent-]{spaces-descent}
\externaldocument[spaces-perfect-]{spaces-perfect}
\externaldocument[spaces-more-morphisms-]{spaces-more-morphisms}
\externaldocument[spaces-flat-]{spaces-flat}
\externaldocument[spaces-groupoids-]{spaces-groupoids}
\externaldocument[spaces-more-groupoids-]{spaces-more-groupoids}
\externaldocument[bootstrap-]{bootstrap}
\externaldocument[spaces-pushouts-]{spaces-pushouts}
\externaldocument[spaces-chow-]{spaces-chow}
\externaldocument[groupoids-quotients-]{groupoids-quotients}
\externaldocument[spaces-more-cohomology-]{spaces-more-cohomology}
\externaldocument[spaces-simplicial-]{spaces-simplicial}
\externaldocument[spaces-duality-]{spaces-duality}
\externaldocument[formal-spaces-]{formal-spaces}
\externaldocument[restricted-]{restricted}
\externaldocument[spaces-resolve-]{spaces-resolve}
\externaldocument[formal-defos-]{formal-defos}
\externaldocument[defos-]{defos}
\externaldocument[cotangent-]{cotangent}
\externaldocument[examples-defos-]{examples-defos}
\externaldocument[algebraic-]{algebraic}
\externaldocument[examples-stacks-]{examples-stacks}
\externaldocument[stacks-sheaves-]{stacks-sheaves}
\externaldocument[criteria-]{criteria}
\externaldocument[artin-]{artin}
\externaldocument[quot-]{quot}
\externaldocument[stacks-properties-]{stacks-properties}
\externaldocument[stacks-morphisms-]{stacks-morphisms}
\externaldocument[stacks-limits-]{stacks-limits}
\externaldocument[stacks-cohomology-]{stacks-cohomology}
\externaldocument[stacks-perfect-]{stacks-perfect}
\externaldocument[stacks-introduction-]{stacks-introduction}
\externaldocument[stacks-more-morphisms-]{stacks-more-morphisms}
\externaldocument[stacks-geometry-]{stacks-geometry}
\externaldocument[moduli-]{moduli}
\externaldocument[moduli-curves-]{moduli-curves}
\externaldocument[examples-]{examples}
\externaldocument[exercises-]{exercises}
\externaldocument[guide-]{guide}
\externaldocument[desirables-]{desirables}
\externaldocument[coding-]{coding}
\externaldocument[obsolete-]{obsolete}
\externaldocument[fdl-]{fdl}
\externaldocument[index-]{index}

% Theorem environments.
%
\theoremstyle{plain}
\newtheorem{theorem}[subsection]{Theorem}
\newtheorem{proposition}[subsection]{Proposition}
\newtheorem{lemma}[subsection]{Lemma}

\theoremstyle{definition}
\newtheorem{definition}[subsection]{Definition}
\newtheorem{example}[subsection]{Example}
\newtheorem{exercise}[subsection]{Exercise}
\newtheorem{situation}[subsection]{Situation}

\theoremstyle{remark}
\newtheorem{remark}[subsection]{Remark}
\newtheorem{remarks}[subsection]{Remarks}

\numberwithin{equation}{subsection}

% Macros
%
\def\lim{\mathop{\mathrm{lim}}\nolimits}
\def\colim{\mathop{\mathrm{colim}}\nolimits}
\def\Spec{\mathop{\mathrm{Spec}}}
\def\Hom{\mathop{\mathrm{Hom}}\nolimits}
\def\Ext{\mathop{\mathrm{Ext}}\nolimits}
\def\SheafHom{\mathop{\mathcal{H}\!\mathit{om}}\nolimits}
\def\SheafExt{\mathop{\mathcal{E}\!\mathit{xt}}\nolimits}
\def\Sch{\mathit{Sch}}
\def\Mor{\mathop{\mathrm{Mor}}\nolimits}
\def\Ob{\mathop{\mathrm{Ob}}\nolimits}
\def\Sh{\mathop{\mathit{Sh}}\nolimits}
\def\NL{\mathop{N\!L}\nolimits}
\def\CH{\mathop{\mathrm{CH}}\nolimits}
\def\proetale{{pro\text{-}\acute{e}tale}}
\def\etale{{\acute{e}tale}}
\def\QCoh{\mathit{QCoh}}
\def\Ker{\mathop{\mathrm{Ker}}}
\def\Im{\mathop{\mathrm{Im}}}
\def\Coker{\mathop{\mathrm{Coker}}}
\def\Coim{\mathop{\mathrm{Coim}}}

% Boxtimes
%
\DeclareMathSymbol{\boxtimes}{\mathbin}{AMSa}{"02}

%
% Macros for moduli stacks/spaces
%
\def\QCohstack{\mathcal{QC}\!\mathit{oh}}
\def\Cohstack{\mathcal{C}\!\mathit{oh}}
\def\Spacesstack{\mathcal{S}\!\mathit{paces}}
\def\Quotfunctor{\mathrm{Quot}}
\def\Hilbfunctor{\mathrm{Hilb}}
\def\Curvesstack{\mathcal{C}\!\mathit{urves}}
\def\Polarizedstack{\mathcal{P}\!\mathit{olarized}}
\def\Complexesstack{\mathcal{C}\!\mathit{omplexes}}
% \Pic is the operator that assigns to X its picard group, usage \Pic(X)
% \Picardstack_{X/B} denotes the Picard stack of X over B
% \Picardfunctor_{X/B} denotes the Picard functor of X over B
\def\Pic{\mathop{\mathrm{Pic}}\nolimits}
\def\Picardstack{\mathcal{P}\!\mathit{ic}}
\def\Picardfunctor{\mathrm{Pic}}
\def\Deformationcategory{\mathcal{D}\!\mathit{ef}}
\usepackage{cleveref}
\crefname{originalreference}{original reference}{original references}
\Crefname{originalreference}{Original reference}{Original references}
\makeatletter
\newlabel{categories-lemma-2-product-categories-over-C}{{0040}{source}{}{sourcepointer-1}{}}
\newlabel{categories-lemma-2-product-categories-over-C@cref}{{[originalreference][0][]0040}{[1][0][]source}{}{}{}}
\makeatother
\crefname{originaldefinition}{original definition}{original definitions}
\Crefname{originaldefinition}{Original definition}{Original definitions}
\makeatletter
\newlabel{definition-algebraic-stack}{{026O}{source}{}{sourcepointer-2}{}}
\newlabel{definition-algebraic-stack@cref}{{[originaldefinition][0][]026O}{[1][0][]source}{}{}{}}
\makeatother
\crefname{originallemma}{original lemma}{original lemmas}
\Crefname{originallemma}{Original lemma}{Original lemmas}
\makeatletter
\newlabel{lemma-base-change-representable-transformations-property}{{045C}{source}{}{sourcepointer-3}{}}
\newlabel{lemma-base-change-representable-transformations-property@cref}{{[originallemma][0][]045C}{[1][0][]source}{}{}{}}
\makeatother
\crefname{originallemma}{original lemma}{original lemmas}
\Crefname{originallemma}{Original lemma}{Original lemmas}
\makeatletter
\newlabel{lemma-composition-representable-transformations-property}{{045B}{source}{}{sourcepointer-4}{}}
\newlabel{lemma-composition-representable-transformations-property@cref}{{[originallemma][0][]045B}{[1][0][]source}{}{}{}}
\makeatother
\crefname{originallemma}{original lemma}{original lemmas}
\Crefname{originallemma}{Original lemma}{Original lemmas}
\makeatletter
\newlabel{lemma-representable-diagonal}{{045G}{source}{}{sourcepointer-5}{}}
\newlabel{lemma-representable-diagonal@cref}{{[originallemma][0][]045G}{[1][0][]source}{}{}{}}
\makeatother
\crefname{originalreference}{original reference}{original references}
\Crefname{originalreference}{Original reference}{Original references}
\makeatletter
\newlabel{stacks-lemma-2-product-stacks-in-groupoids}{{02ZL}{source}{}{sourcepointer-6}{}}
\newlabel{stacks-lemma-2-product-stacks-in-groupoids@cref}{{[originalreference][0][]02ZL}{[1][0][]source}{}{}{}}
\makeatother
\title{2-fibre products of algebraic stacks over a stack with space-representable diagonal are algebraic}
\author{The Stacks Project Authors}
\date{stacks-project:algebraic.tex:local-snapshot-20261012; GNU FDL 1.2 or later}
\begin{document}
\maketitle
\noindent\textit{Editorial note.} No omission markers in-span; the Isom square commutativity is read off the explicit object descriptions (mechanical). The two property-transfer inputs are not carded ('Proof omitted.' each) and are declared external.\par\medskip
\noindent\textit{Allowed standard prerequisites (editorial).}\begin{itemize}
\item Stack closure under 2-fibre products. Stacks, lemma 2-product-stacks-in-groupoids (cross-chapter, named). Condition (1).
\item Diagonal criterion. lemma-representable-diagonal (carded L26680). Isom sheaves of algebraic-stack inputs are algebraic spaces.
\item Base-change property transfer. lemma-base-change-representable-transformations-property (same chapter, NOT carded, declared external). Properties transfer along base change.
\item Composition property transfer. lemma-composition-representable-transformations-property (same chapter, NOT carded, declared external). Properties transfer along composition.
\end{itemize}\medskip
\ifdefined\hypertarget\hypertarget{sourcepointer-1}{}\fi
\noindent\textit{Source reference.} Original reference 0040: focused-seeds/geometry/S9-alg-2-fibre-product-general/ref-tags.txt, source lines 5--5. Navigation pointer only: the counted proof cites this statement by \textbackslash{}ref as a standing input; the deduction the counted proof has to make is carried out inside the counted span, and no part of that deduction is delegated to this pointer.\par\medskip
\ifdefined\hypertarget\hypertarget{sourcepointer-2}{}\fi
\noindent\textit{Source reference.} Original definition 026O: sources/stacks-algebraic/pinned-original-source.tex, source lines 1506--1506. Navigation pointer only: the counted proof cites this statement by \textbackslash{}ref as a standing input; the deduction the counted proof has to make is carried out inside the counted span, and no part of that deduction is delegated to this pointer.\par\medskip
\ifdefined\hypertarget\hypertarget{sourcepointer-3}{}\fi
\noindent\textit{Source reference.} Original lemma 045C: sources/stacks-algebraic/pinned-original-source.tex, source lines 1180--1180. Navigation pointer only: the counted proof cites this statement by \textbackslash{}ref as a standing input; the deduction the counted proof has to make is carried out inside the counted span, and no part of that deduction is delegated to this pointer.\par\medskip
\ifdefined\hypertarget\hypertarget{sourcepointer-4}{}\fi
\noindent\textit{Source reference.} Original lemma 045B: sources/stacks-algebraic/pinned-original-source.tex, source lines 1159--1159. Navigation pointer only: the counted proof cites this statement by \textbackslash{}ref as a standing input; the deduction the counted proof has to make is carried out inside the counted span, and no part of that deduction is delegated to this pointer.\par\medskip
\ifdefined\hypertarget\hypertarget{sourcepointer-5}{}\fi
\noindent\textit{Source reference.} Original lemma 045G: sources/stacks-algebraic/pinned-original-source.tex, source lines 1337--1337. Navigation pointer only: the counted proof cites this statement by \textbackslash{}ref as a standing input; the deduction the counted proof has to make is carried out inside the counted span, and no part of that deduction is delegated to this pointer.\par\medskip
\ifdefined\hypertarget\hypertarget{sourcepointer-6}{}\fi
\noindent\textit{Source reference.} Original reference 02ZL: focused-seeds/geometry/S9-alg-2-fibre-product-general/ref-tags.txt, source lines 6--6. Navigation pointer only: the counted proof cites this statement by \textbackslash{}ref as a standing input; the deduction the counted proof has to make is carried out inside the counted span, and no part of that deduction is delegated to this pointer.\par\medskip
\begin{lemma}
\label{lemma-2-fibre-product-general}
Let $S$ be a scheme contained in $\Sch_{fppf}$.
Let $\mathcal{Z}$ be a stack in groupoids over $(\Sch/S)_{fppf}$
whose diagonal is representable by algebraic spaces.
Let $\mathcal{X}$, $\mathcal{Y}$ be algebraic stacks over $S$.
Let $f : \mathcal{X} \to \mathcal{Z}$, $g : \mathcal{Y} \to \mathcal{Z}$
be $1$-morphisms of stacks in groupoids. Then the $2$-fibre product
$\mathcal{X} \times_{f, \mathcal{Z}, g} \mathcal{Y}$ is an algebraic stack.
\end{lemma}

\begin{proof}
We have to check conditions (1), (2), and (3) of
Definition \ref{definition-algebraic-stack}.
The first condition follows from
Stacks, Lemma \ref{stacks-lemma-2-product-stacks-in-groupoids}.

\medskip\noindent
The second condition we have to check is that the $\mathit{Isom}$-sheaves
are representable by algebraic spaces. To do this, suppose that
$T$ is a scheme over $S$, and $u, v$ are objects of
$(\mathcal{X} \times_{f, \mathcal{Z}, g} \mathcal{Y})_T$.
By our construction of $2$-fibre products (which goes all the way
back to
Categories, Lemma \ref{categories-lemma-2-product-categories-over-C})
we may write $u = (x, y, \alpha)$ and $v = (x', y', \alpha')$.
Here $\alpha : f(x) \to g(y)$ and similarly for $\alpha'$.
Then it is clear that
$$
\xymatrix{
\mathit{Isom}(u, v) \ar[d] \ar[rr] & &
\mathit{Isom}(y, y') \ar[d]^{\phi \mapsto g(\phi) \circ \alpha} \\
\mathit{Isom}(x, x') \ar[rr]^-{\psi \mapsto \alpha' \circ f(\psi)} & &
\mathit{Isom}(f(x), g(y'))
}
$$
is a cartesian diagram of sheaves on $(\Sch/T)_{fppf}$.
Since by assumption the sheaves
$\mathit{Isom}(y, y')$, $\mathit{Isom}(x, x')$, $\mathit{Isom}(f(x), g(y'))$
are algebraic spaces (see
Lemma \ref{lemma-representable-diagonal})
we see that $\mathit{Isom}(u, v)$
is an algebraic space.

\medskip\noindent
Let $U, V \in \Ob((\Sch/S)_{fppf})$,
and let $x, y$ be surjective smooth morphisms
$x : (\Sch/U)_{fppf} \to \mathcal{X}$,
$y : (\Sch/V)_{fppf} \to \mathcal{Y}$.
Consider the morphism
$$
(\Sch/U)_{fppf}
\times_{f \circ x, \mathcal{Z}, g \circ y}
(\Sch/V)_{fppf}
\longrightarrow
\mathcal{X} \times_{f, \mathcal{Z}, g} \mathcal{Y}.
$$
As the diagonal of $\mathcal{Z}$ is representable by algebraic spaces
the source of this arrow is representable by an algebraic space $F$, see
Lemma \ref{lemma-representable-diagonal}.
Moreover, the morphism is the composition
of base changes of $x$ and $y$, hence surjective and smooth, see
Lemmas \ref{lemma-base-change-representable-transformations-property} and
\ref{lemma-composition-representable-transformations-property}.
Choosing a scheme $W$ and a surjective \'etale morphism $W \to F$
we see that the composition of the displayed $1$-morphism
with the corresponding $1$-morphism
$$
(\Sch/W)_{fppf}
\longrightarrow
(\Sch/U)_{fppf}
\times_{f \circ x, \mathcal{Z}, g \circ y}
(\Sch/V)_{fppf}
$$
is surjective and smooth which proves the last condition.
\end{proof}

\end{document}
